\documentclass[11pt]{article}
\usepackage[T1]{fontenc}
\usepackage{lmodern}
\usepackage[margin=1in]{geometry}
\usepackage{amsmath,amssymb,amsthm,mathtools}
\usepackage{microtype}
\usepackage[colorlinks=true,linkcolor=blue,citecolor=blue,urlcolor=blue]{hyperref}
\newtheorem{theorem}{Theorem}
\newtheorem{lemma}[theorem]{Lemma}
\newtheorem{proposition}[theorem]{Proposition}
\newcommand{\Li}{\operatorname{Li}}
\newcommand{\dd}{\,\mathrm{d}}
\title{A parameter-derivative proof of Sun's\\
binomial-harmonic identity for $\zeta(5)$}
\author{GPT-6 Astra}
\date{November 1, 2026}
\begin{document}
\maketitle
\begin{abstract}
We prove Conjecture 3.11, equation (3.15), of Zhi-Wei Sun's
\emph{Various conjectural series identities}. The series has the value
$\frac{93}{2}\zeta(5)+\pi^4\log 2$.
The proof differentiates a parameter identity already established by
Kam Cheong Au and uses a second evaluation proved in the same work.
We give the normalization, justify differentiation by locally uniform
convergence, and evaluate the remaining harmonic sums explicitly.
We also correct the Gamma interpolation appearing in Sun's accompanying
remark. Exact rational interval computations provide an independent check
of the final identity.
\end{abstract}

\section{Statement and relation to earlier work}
Write $H_0=0$, $H_m=\sum_{j=1}^m j^{-1}$ for $m\geq1$, and
$\zeta(s)=\sum_{j=1}^{\infty}j^{-s}$ for real $s>1$.
All logarithms below are natural. Define
\begin{equation}\label{eq:definitions}
 P(x)=21x^3-22x^2+8x-1,\qquad
 c_n=\frac{256^n}{n^7\binom{2n}{n}^7},\qquad
 h_n=H_{2n-1}-H_{n-1}.
\end{equation}
\begin{theorem}\label{thm:main}
The following series converges absolutely and satisfies
\begin{equation}\label{eq:target}
 \boxed{\displaystyle
 \sum_{n=1}^{\infty}
 \frac{256^n\bigl(14P(n)(H_{2n-1}-H_{n-1})-63n^2+44n-8\bigr)}
 {n^7\binom{2n}{n}^7}
 =\frac{93}{2}\zeta(5)+\pi^4\log2.}
\end{equation}
\end{theorem}

This is exactly Sun's Conjecture 3.11, equation (3.15), in
\cite{Sun}, version 3, dated April 13, 2026. The conjecture itself is
dated April 12, 2026. Its companion identity
\begin{equation}\label{eq:companion}
 \sum_{n=1}^{\infty}c_nP(n)=\frac{\pi^4}{8}
\end{equation}
was conjectured by Guillera and proved by Au in
\cite[Example XI, equation (4.2)]{Au}.

Our two substantive hypergeometric inputs are Au's parameter identity
\cite[equation (4.4)]{Au} and the evaluation at the end of
\cite[Example XV]{Au}. Both are stated explicitly in the next section.
Au derives them using Wilf--Zeilberger pairs; the first combines two such
pairs, including their boundary contributions. The result here is the
explicit derivative consequence giving Sun's later conjecture.
A source and subsequent-proof search performed on November 1, 2026
is recorded in the accompanying literature review.

\section{Two established hypergeometric inputs}
For $m\geq0$, let $(z)_m=\prod_{j=0}^{m-1}(z+j)$, with $(z)_0=1$.
The Gamma function is $\Gamma(z)=\int_0^{\infty}t^{z-1}e^{-t}\dd t$
for $\Re z>0$, continued meromorphically elsewhere. Its logarithmic
derivative is denoted by $\psi$.

For $a$ near zero define
\begin{align}
 S(a)&=\sum_{n=1}^{\infty}
 \frac{64^{-n}P(n+2a)(1+2a)_n^7}
 {(n+2a)^7(\tfrac12+2a)_n^7},\label{eq:S}\\
 C(a)&=\frac{\Gamma(\tfrac12-a)\Gamma(\tfrac12+2a)^4\Gamma(1+2a)}
 {\Gamma(\tfrac12+a)\Gamma(1+a)^4\Gamma(1+4a)},\label{eq:C}\\
 B(a)&=\sum_{k=0}^{\infty}
 \frac{(3a+2k+1)(\tfrac12+a)_k^3(1+a)_k}
 {(4a+2k+1)^3(\tfrac12+2a)_k^3(1+2a)_k},\label{eq:B}\\
 A(a)&=\sum_{k=0}^{\infty}
 \frac{(6a+4k+3)(\tfrac12+a)_k(1+a)_k^4(1+4a)_k}
 {(2k+1-2a)(k+1+2a)(4a+2k+1)^4
 (\tfrac12-a)_k(\tfrac12+2a)_k^4(1+2a)_k}.
 \label{eq:A}
\end{align}

\begin{proposition}[Au's identities in the normalization used here]
For $a$ in a neighborhood of zero,
\begin{equation}\label{eq:Au}
 S(a)=C(a)B(a)-4aA(a).
\end{equation}
In addition,
\begin{equation}\label{eq:AuXV}
 \sum_{n=1}^{\infty}
 \frac{(1)_n^4}{(\tfrac12)_n^4}
 \frac{4n-1}{4n^5(2n-1)}
 =93\zeta(5)-7\pi^2\zeta(3).
\end{equation}
\end{proposition}

These are established results of \cite{Au}.
To check the translation of equation (4.4), its first series is
$32A(a)$, its second is $-8C(a)B(a)/a$, and its right side is
$-8S(a)/a$. The last assertion follows from
\begin{equation}\label{eq:dup}
 (1+4a)_{2n}=4^n(\tfrac12+2a)_n(1+2a)_n
\end{equation}
and expansion of $P(n+2a)$. Multiplication by $-a/8$ gives
\eqref{eq:Au}. Formula \eqref{eq:AuXV} is the final displayed summation
of Au's Example XV, before Remark 5.3. In that example, it is obtained
from equation (5.2) by the degree-two coefficient combination specified
in footnote 21.

\begin{lemma}\label{lem:analytic}
The functions $A$, $B$, and $S$ are holomorphic for $|a|<1/8$ and may
be differentiated term by term there. Equation \eqref{eq:Au} extends
holomorphically through $a=0$.
\end{lemma}
\begin{proof}
Fix $0<r<1/8$. All summands are holomorphic on a neighborhood of
$|a|\leq r$; their denominators are nonzero there. The uniform Gamma-ratio
estimate
\[
 \frac{\Gamma(k+u)}{\Gamma(k+v)}
 =k^{u-v}\bigl(1+O(k^{-1})\bigr)
\]
for $u,v$ in compact sets shows that the summands of $A$, $B$, and $S$
are respectively
\[
 O(k^{-3}),\qquad O(k^{-2+4r}),\qquad
 O(64^{-n}n^{-1/2}),
\]
uniformly on that disk. For $A$, the Pochhammer quotient has exponent
$2$ and the rational factor has exponent $-5$. For $B$, these exponents
are $-4a$ and $-2$. For $S$, they are $7/2$ and $-4$.
The finitely many initial terms cause no problem. The three majorants
are summable. Normal convergence and the holomorphic convergence theorem
give the stated differentiability, also by applying Cauchy's formula
on a slightly larger disk. Both sides of \eqref{eq:Au} are therefore
holomorphic at zero, which removes the apparent pole in Au's original
normalization. This argument supplies the continuation and
differentiation needed below.
\end{proof}

\section{The harmonic sums produced by differentiation}
\begin{lemma}\label{lem:Euler}
One has
\begin{align}
 U:=\sum_{k=0}^{\infty}\frac{H_k}{(2k+1)^2}
 &=\frac74\zeta(3)-\frac{\pi^2}{4}\log2,\label{eq:U}\\
 W:=\sum_{k=0}^{\infty}\frac{H_{2k}}{(2k+1)^2}
 &=\frac7{16}\zeta(3).\label{eq:W}
\end{align}
\end{lemma}
\begin{proof}
The generating function
$\sum_{k\geq0}H_kz^k=-\log(1-z)/(1-z)$ gives
\[
 U=\int_0^1\frac{\log x\log(1-x^2)}{1-x^2}\dd x
 =\frac14\int_0^1
 \frac{y^{-1/2}\log y\log(1-y)}{1-y}\dd y.
\]
For $\alpha>0$, expand the beta function at $b=0$:
\begin{equation}\label{eq:beta}
 \mathrm B(\alpha,b)=\frac1b-\gamma-\psi(\alpha)
 +\frac b2\bigl((\gamma+\psi(\alpha))^2+\zeta(2)-\psi'(\alpha)\bigr)
 +O(b^2).
\end{equation}
Here $\gamma$ is Euler's constant; the expansion follows by multiplying
the Taylor series of $\Gamma(b)\Gamma(\alpha)/\Gamma(\alpha+b)$.
The derivative in $\alpha$ removes the pole. Taking the mixed derivative
at $\alpha=1/2$, $b=0$ gives
\[
 U=\frac14\bigl((\gamma+\psi(\tfrac12))\psi'(\tfrac12)
 -\tfrac12\psi''(\tfrac12)\bigr).
\]
Differentiation under the integral is valid near these parameters:
at $y=1$ the factor $\log y$ cancels the apparent singularity, leaving
an integrable logarithm. The values
\[
 \psi(\tfrac12)=-\gamma-2\log2,\qquad
 \psi'(\tfrac12)=\frac{\pi^2}{2},\qquad
 \psi''(\tfrac12)=-14\zeta(3)
\]
follow from Gamma duplication and the absolutely convergent polygamma
series. They prove \eqref{eq:U}.

Taking the even part of the harmonic generating function yields
\begin{equation}\label{eq:Wint}
 W=\frac12\int_0^1\log x
 \left(\frac{\log(1-x)}{1-x}+\frac{\log(1+x)}{1+x}\right)\dd x.
\end{equation}
Integration by parts and then the substitution $t=1-x$ give
\[
 \int_0^1\frac{\log x\log(1-x)}{1-x}\dd x
 =\frac12\int_0^1\frac{\log^2(1-x)}x\dd x
 =\frac12\int_0^1\frac{\log^2t}{1-t}\dd t=\zeta(3).
\]
The second integral in \eqref{eq:Wint} is $-J/2$, where
$J=\int_0^1\log^2(1+x)\dd x/x$. We give the details of its evaluation.

Define $\Li_s(z)=\sum_{n\geq1}z^n/n^s$ on the unit disk, using continuous
endpoint values where required. Differentiating and fixing constants
at zero proves
\begin{align}
 \Li_2(x)+\Li_2(1-x)&=\zeta(2)-\log x\log(1-x),\label{eq:dilog1}\\
 \Li_2\!\left(-\frac{x}{1-x}\right)
 &=-\Li_2(x)-\frac12\log^2(1-x).\label{eq:dilog2}
\end{align}
For $0<x\leq1/2$, these in turn prove by differentiation
\begin{align}
 &\Li_3(x)+\Li_3(1-x)+\Li_3\!\left(-\frac{x}{1-x}\right)\notag\\
 &\hspace{1em}=\zeta(3)+\zeta(2)\log(1-x)
 -\frac12\log x\log^2(1-x)+\frac16\log^3(1-x).
 \label{eq:trilog}
\end{align}
Writing $L=\log2$, substituting $x=1/2$, and using
$\Li_3(-1)=-3\zeta(3)/4$, we obtain
\[
 \Li_2(\tfrac12)=\frac{\pi^2}{12}-\frac{L^2}{2},\qquad
 \Li_3(\tfrac12)=\frac78\zeta(3)-\frac{\pi^2L}{12}+\frac{L^3}{6}.
\]
In $J$, the substitution $t=(1+x)^{-1}$ gives
\[
 J=\int_{1/2}^1\frac{\log^2t}{t(1-t)}\dd t
 =\frac{L^3}{3}+2\zeta(3)-\int_0^{1/2}\frac{\log^2t}{1-t}\dd t.
\]
An antiderivative of the last integrand is
$-\log^2t\log(1-t)-2\log t\Li_2(t)+2\Li_3(t)$.
Consequently
\[
 J=2\zeta(3)-\frac23L^3-2L\Li_2(\tfrac12)
 -2\Li_3(\tfrac12)=\frac14\zeta(3).
\]
Equation \eqref{eq:Wint} now gives $W=7\zeta(3)/16$.
All integrals used here converge absolutely; the generating-function
integrations also follow from monotone convergence after adjusting signs.
\end{proof}

\section{Proof of the conjecture}
The values in \eqref{eq:C} and \eqref{eq:B} at zero are
\begin{equation}\label{eq:CB0}
 C(0)=\pi^2,\qquad C'(0)=-12\pi^2\log2,\qquad B(0)=\frac{\pi^2}{8}.
\end{equation}
Indeed, the logarithmic derivative of $C$ is
$6\psi(1/2)-6\psi(1)=-12\log2$ at zero.
The summand of $B(0)$ is $(2k+1)^{-2}$.
Logarithmic differentiation of each factor of \eqref{eq:B} gives
\begin{equation}\label{eq:Bprime}
 B'(0)=\sum_{k=0}^{\infty}
 \left(\frac{-6H_{2k}+2H_k}{(2k+1)^2}
 -\frac9{(2k+1)^3}\right)
 =-7\zeta(3)-\frac{\pi^2}{2}\log2.
\end{equation}
For clarity, the Pochhammer factors contribute $-6H_{2k}+2H_k$,
using $\sum_{j=0}^{k-1}(j+1/2)^{-1}=2H_{2k}-H_k$.
The linear numerator and cubed linear denominator contribute
$-9/(2k+1)$. The last equality in \eqref{eq:Bprime} follows from
Lemma~\ref{lem:Euler} and
$\sum_{k\geq0}(2k+1)^{-3}=7\zeta(3)/8$.

Setting $a=0$ in \eqref{eq:A} gives
\[
 A(0)=\sum_{k=0}^{\infty}
 \frac{(4k+3)(1)_k^4}{(2k+1)^5(k+1)(\tfrac12)_k^4}.
\]
With $n=k+1$ and
$\frac{(1)_{n-1}}{(1/2)_{n-1}}
=\frac{n-1/2}{n}\frac{(1)_n}{(1/2)_n}$, its summand is one quarter
of the corresponding summand in \eqref{eq:AuXV}. Hence
\begin{equation}\label{eq:A0}
 A(0)=\frac{93\zeta(5)-7\pi^2\zeta(3)}4.
\end{equation}
Differentiating \eqref{eq:Au}, justified by Lemma~\ref{lem:analytic},
and applying \eqref{eq:CB0}--\eqref{eq:A0}, we find
\begin{align}
 S'(0)&=C'(0)B(0)+C(0)B'(0)-4A(0)\notag\\
 &=-\frac32\pi^4\log2-7\pi^2\zeta(3)
 -\frac12\pi^4\log2-93\zeta(5)+7\pi^2\zeta(3)\notag\\
 &=-93\zeta(5)-2\pi^4\log2.\label{eq:Sprime}
\end{align}

To identify this derivative with the requested series, set
\[
 q_n(a)=\frac{64^{-n}(1+2a)_n^7}
 {(n+2a)^7(\tfrac12+2a)_n^7}.
\]
The elementary identity $(1/2)_n=(2n)!/(4^nn!)$ gives $q_n(0)=c_n$.
Differentiating the finite products directly gives
\begin{align*}
 \frac{q_n'(0)}{q_n(0)}
 &=14H_n-14(2H_{2n}-H_n)-\frac{14}{n}\\
 &=-28(H_{2n-1}-H_{n-1})=-28h_n.
\end{align*}
Since $P'(n)=63n^2-44n+8$, it follows that
\[
 S'(0)=\sum_{n=1}^{\infty}c_n\bigl(2P'(n)-28P(n)h_n\bigr)
 =-2\sum_{n=1}^{\infty}c_n\bigl(14P(n)h_n-P'(n)\bigr).
\]
Combining this with \eqref{eq:Sprime} proves \eqref{eq:target}.
Absolute convergence follows either from the same exponential
majorant or from the explicit bound in Section~\ref{sec:check}.
\hfill$\square$

\section{The Gamma interpolation in the source}
Sun's Remark 3.11 gives an auxiliary Gamma expression accompanying
the conjecture. The denominator printed there has $\Gamma(2x)$;
the expression that interpolates the stated binomial summand is
\begin{equation}\label{eq:corrected}
 f(x)=\frac{256^x P(x)\Gamma(x)^{14}}{128\Gamma(2x)^7}
 =\pi^{7/2}64^{-x}P(x)\frac{\Gamma(x)^7}{\Gamma(x+1/2)^7},
 \qquad x>0.
\end{equation}
The factor $\Gamma(2x)^7$ follows directly by replacing factorials
in \eqref{eq:definitions}; Gamma duplication proves the second equality.
In particular $f(n)=c_nP(n)$. If the seventh power is omitted,
this equality already fails at $n=2$, by a factor $\Gamma(4)^6=6^6$.

Logarithmic differentiation of the weight in \eqref{eq:corrected} gives
\[
 f'(n)=c_n\bigl(P'(n)+(8\log2-14h_n)P(n)\bigr).
\]
The companion identity \eqref{eq:companion} also follows by putting
$a=0$ in \eqref{eq:Au}. Thus the theorem gives the equivalent
derivative evaluation
\[
 \sum_{n=1}^{\infty}f'(n)=-\frac{93}{2}\zeta(5).
\]
The proof of the theorem used the finite-product derivative directly;
the corrected interpolation provides a separate check of the formulation.

\section{Independent arithmetic checks and reproducibility}\label{sec:check}
The proof above is analytic and uses the two explicitly attributed
identities of Au. The included verification program checks the algebraic
normalization and constant cancellation using exact symbolic arithmetic,
and computes independent rational enclosures of both sides of
\eqref{eq:target}. We specify its error bounds here.

Since $0<h_n\leq1$, the absolute value of the numerator in parentheses
in \eqref{eq:target} is at most $843n^3$ for $n\geq1$.
The largest coefficient in $(1+1)^{2n}$ is $\binom{2n}{n}$, so
$\binom{2n}{n}\geq4^n/(2n+1)$. Therefore the absolute remainder
after $N$ terms is bounded by
\begin{equation}\label{eq:tail}
 R_N\leq843\cdot3^7\sum_{n>N}\frac{n^3}{64^n}
 \leq\frac87\,843\cdot3^7\frac{(N+1)^3}{64^{N+1}}.
\end{equation}
The last step uses the successive-term ratio at most $1/8$.
For $N=100$ this is less than $8.175\times10^{-171}$.

For the other side, the program uses Machin's formula
$\pi=16\arctan(1/5)-4\arctan(1/239)$ with alternating-series bounds,
and the positive series
$\log2=2\sum_{j\geq0}3^{-2j-1}/(2j+1)$ with a geometric tail bound.
For $\zeta(5)$, Euler's transformation gives
\[
 \frac{15}{16}\zeta(5)=\sum_{j=0}^{\infty}\frac{d_j}{2^{j+1}},
 \qquad
 d_j=\sum_{i=0}^{j}(-1)^i\binom ji\frac1{(i+1)^5}.
\]
The integral
$d_j=\frac1{24}\int_0^1(1-t)^j(-\log t)^4\dd t$
shows $0<d_j\leq1$. The omitted tail after $M$ terms is therefore
at most $2^{-M}$. In particular, the rational computations require
no floating-point error model.

Both enclosures round outward to the same interval of width $10^{-160}$.
For a readable numerical example their common value begins
\[
 115.7359774252925386770599308882310257130715848179365598448183046\ldots.
\]
The first term of the series is exactly $114$.
The arithmetic calculation is a consistency check accompanying the proof.

To reproduce the checks, install \path{requirements.txt} and run
\texttt{python src/verify.py} from the supplement root. It writes
\path{certificates/verification.json}. The source \path{paper/main.tex}
builds with pdfLaTeX; exact build commands are in the README.

\paragraph{Authorship.}
This manuscript was generated by the AI system GPT-6 Astra.
The accompanying scripts check the stated algebra and numerical bounds.

\begin{thebibliography}{9}
\bibitem{Sun}
Zhi-Wei Sun,
\emph{Various conjectural series identities},
arXiv:2603.29973v3, April 13, 2026.
Conjecture 3.11, equation (3.15), and Remark 3.11.
\url{https://arxiv.org/abs/2603.29973v3}.

\bibitem{Au}
Kam Cheong Au,
\emph{Wilf--Zeilberger seeds and non-trivial hypergeometric identities},
Journal of Symbolic Computation \textbf{130} (2025), 102421.
\href{https://doi.org/10.1016/j.jsc.2025.102421}{doi:10.1016/j.jsc.2025.102421}.
The version used for equation numbering is arXiv:2312.14051v4,
October 10, 2024, Examples XI and XV.
\url{https://arxiv.org/abs/2312.14051v4}.
\end{thebibliography}
\end{document}
